\chapter{Conclusions and Future Work} \label{ch:conclusions}

This project investigated whether a biologically inspired neuromorphic model could perform three-dimensional active acoustic localisation in a way that was both functional and interpretable. The final system estimates target distance, azimuth, and elevation from a simulated echolocation call using a cochlear spike front end, separate distance/azimuth/elevation pathway models, SC-style population readouts, and a small trainable SNN fusion stage.

The main conclusion is that biological structure was useful, but not sufficient on its own. The best performance was not obtained by training a small SNN directly from the waveform, nor by relying only on the fixed hand-designed pathway outputs. Instead, the strongest architecture used the biomimetic pathways as structured feature extraction, then added a residual SNN readout to learn context-dependent corrections. The clean raw pathway COM already achieved a combined error of 0.0585, outperforming the strongest clean input-only SNN baseline at 0.0828 before residual training was introduced. This supports the central argument of the report, that simplified biological pathways can provide useful intermediate representations for neuromorphic localisation, while limited training can improve final integration without removing interpretability.

\section{Summary of Contributions}

The first contribution was a complete simulated active-acoustic localisation pipeline. A frequency-modulated call was emitted, reflected from a target, and returned as a binaural echo containing distance delay, attenuation, binaural azimuth cues, and elevation-dependent spectral shaping. This provided a controlled environment in which the physical cues for distance, azimuth, and elevation could be tested separately before being combined.

The second contribution was a set of interpretable neuromorphic pathway models. The cochlear front end produced a frequency-organised spike representation, the distance pathway used pulse-echo coincidence, the azimuth pathway used ITD coincidence and ILD opponent comparison, and the elevation pathway used a DCN-like spectral-template model for comb-filter cues. These pathway models were simplified, but each was tied to a clear acoustic cue and a plausible biological computation.

The third contribution was the comparison of fixed and trainable readouts. Direct centre-of-mass decoding, calibrated decoding, and SC-style continuous-attractor readouts were tested as population-code readouts. The final residual SNN then used the raw pathway populations, fixed readout estimates, and confidence features to improve the coordinate estimate. In the constrained noise-free test region, this residual pathway-feature SNN achieved a combined normalised error of 0.022 and a Euclidean MAE of 0.125~m. Under environmental noise, it reduced the combined error from 0.263 for the fixed CANN/calibrated baseline to 0.055. Under environmental noise plus controlled delayed echoes, it maintained a similar combined error of 0.055.

\section{Main Findings}

\paragraph{Biological structure provided meaningful feature extraction.} Raw waveform SNN baselines performed poorly in the tested regime, cochlear-raster baselines performed better, and the best results were obtained when the trainable SNN received the full pathway-derived feature set. This suggests that the hand-designed pathways performed meaningful feature extraction. Time-of-flight estimation, binaural comparison, and spectral-template matching were represented explicitly, allowing the trainable readout to focus on residual errors and cue interactions.

\paragraph{Distance was the most stable coordinate pathway.} The three coordinate pathways did not fail in the same way. Distance was the most stable in the tuned 0.25--5~m operating region. It was also the only pathway with explicit downstream noise-control mechanisms, including onset consensus, DNLL-style suppression, and FM sweep facilitation. These mechanisms may contribute to its stability, although this remains a hypothesis until they are isolated through causal ablation.

\paragraph{Azimuth required a context-dependent choice of cue.} The calibrated ILD branch reached $0.982^\circ$ MAE in isolation, but increased to $6.175^\circ$ in the clean constrained 3D test because its level cue changed with distance. ITD was less accurate in isolation but more robust in clean full 3D, reaching $3.351^\circ$ MAE. The fixed scalar readout therefore uses ITD, while the fusion SNN receives both populations. This retains complementary information without assuming that either cue is sufficient in every context.

\paragraph{Elevation was the most sensitive pathway.} Elevation depends on spectral shaping rather than a direct timing cue. Environmental noise corrupted this cue before DCN template matching, causing the fixed elevation estimate to collapse. The collapse still retained an ordered trend with true elevation. This helps explain why residual fusion could reduce elevation MAE from $26.642^\circ$ to $3.505^\circ$. The fixed scalar estimate failed, but the wider population representation still contained structured information.

\paragraph{Population codes retained information lost by scalar readouts.} Raw pathway populations contributed more than the three CANN scalar readouts when each group was removed in isolation. A population retains peak width, ambiguity, margins, and confidence information that a single COM estimate discards. The clean residual system also achieved sub-bin precision for distance and elevation. This shows the value of distributed readout, but does not establish biological hyperacuity.

\paragraph{The CANN readout was a stabiliser, not an optimiser.} The attractor readout improved some estimates, made little difference to others, and could slightly worsen a readout when the upstream population was already sharp. It should therefore be treated as a neural population readout and stabiliser that has memory capabilities, not as an optimisation stage that automatically improves accuracy. Its recurrent state still provides a route towards memory and temporal tracking when successive calls are processed over time.

\paragraph{Residual fusion was the strongest final architecture.} It preserves the fixed pathway estimate and only asks the SNN to learn a correction, matching the modelling philosophy of the project to use biological structure where the computation is understood, and use training where manual calibration of cue interactions becomes difficult. The current fusion model receives cached population features repeated across its simulation timesteps, rather than temporally aligned pathway spike trains. It can learn contextual relationships between population shape, confidence, and error, but it cannot yet test whether cross-pathway timing alignment provides additional robustness.

\section{Limitations}

The model remains a prototype. The strongest results were obtained in the constrained region for which the system was tuned: 0.25--5~m and $\pm45^\circ$ azimuth/elevation. This is a meaningful operating region for evaluating the proposed architecture, but it is not a general solution to all acoustic localisation.

The acoustic environment is also simplified. The model uses simulated calls, simplified head-shadow effects, and a comb-filter elevation cue rather than measured bat HRTFs or measured target reflections. The delayed-echo condition is a controlled clutter stress test rather than realistic reverberation: its copies retain target-consistent filtering and occur in the same order on every trial. This was appropriate for controlled model development, but it means the results should be treated as proof-of-principle evidence rather than direct evidence of real-world performance.

Many biological mechanisms are represented functionally rather than biophysically. The cochlea is an engineered IIR filterbank, the VCN/VNLL stages are simplified onset processors, and the DCN, IC, AC, and SC are compact computational abstractions. The model should therefore not be claimed to reproduce bat auditory neurophysiology in detail.

The final limitation is that the system remains largely hand tuned and single-object. Many pathway parameters were selected through iterative testing, with training restricted to the fusion readout. This preserved interpretability, but the resulting model is unlikely to be globally optimal. Mechanisms such as DNLL-like suppression also assume one dominant echo, so cluttered scenes and multiple reflectors would require further development. Low-power operation remains a design motivation rather than a measured result.

\section{Future Work}

The first desire would be to move from single-pulse localisation to continuous active sensing with temporally aligned pathway spikes. A more complete system would retain state across calls, allowing the SC-style population readout to act as a temporal estimator rather than a static decoder. This would test whether cross-pathway timing coincidences provide additional robustness, while also supporting smoothing, rejection of transient errors, velocity estimation, and short-term prediction of moving targets. Preparing for this continuous feedback system was the motivation for the CANN readout mechanism.

Another direction is to improve realism and trainability without discarding interpretability. More realistic acoustic tests should use measured HRTFs, geometry-dependent reflections, larger datasets, and more varied noise and multipath conditions. These tests should include multi-object scenes, where DNLL-style suppression and the assumption of one dominant echo will need to be reconsidered. Within the model, selected parameters could be made trainable while preserving pathway structure, such as cochlear gains, coincidence thresholds, ILD opponent weights, elevation-template weights, CANN gains, or calibration parameters.

Finally, the neuromorphic motivation should be tested directly. Future work should measure spike counts, synaptic operations, memory movement, latency, and deployment cost on embedded or neuromorphic hardware. A biologically inspired SNN is only compelling as an engineering system if its structure leads to useful trade-offs in accuracy, power, latency, or debuggability.

\section{Final Remarks}

The project demonstrates a workable route between two unsatisfactory extremes. A purely hand-designed model is interpretable but struggles with complex cue interactions. A purely trained model can fit data but gives little insight into how localisation is being performed. The best system developed here combines both approaches: fixed biomimetic pathways extract meaningful auditory features, and a small trainable SNN performs contextual fusion and residual correction. Comparable neuromorphic studies rarely evaluate active distance estimation alongside both angular coordinates, so the contribution is the modular 3D pipeline and design method rather than a claim of state-of-the-art performance across incompatible benchmarks.

The final conclusion is therefore not that this model proves neuromorphic sensing is lower power, or that it reproduces a bat brain in detail. The defensible conclusion is more specific: biologically structured neuromorphic modelling provided a practical and interpretable feature-extraction strategy for three-dimensional active acoustic localisation. The term biologically informed SNN (BI-SNN) is used here as a project-level description of this middle ground between detailed biological simulation and black-box training, rather than as an established model class. The model remains a prototype, but it establishes a clear foundation for more realistic, continuous, and hardware-aware neuromorphic echolocation systems.
